\section{Validación cruzada}

\subsection{Diseño}

Después de fijar hiperparámetros, arquitectura o ensemble se unieron
entrenamiento y validación en un conjunto de desarrollo de 28\,418 imágenes.
Sobre él se ejecutó StratifiedKFold de cinco partes, con barajado y semilla 42.
La estratificación utilizó la combinación clase--entorno para impedir que uno
de los folds perdiera los pocos ejemplos de roya de campo o concentrara los
fondos de laboratorio.

En cada fold la cabeza se ajustó desde cero sobre cuatro partes y se evaluó en
la quinta. Cuando la selección fue un ensemble, se reajustó una cabeza para
cada backbone y se conservaron los pesos de voting aprendidos antes. No se
volvieron a optimizar pesos dentro de los folds. Esto mide estabilidad de la
configuración fijada y evita cinco búsquedas nuevas.

\begin{table}[H]
  \centering
  \caption{Resultados de cinco folds, media y desviación estándar muestral.}
  \label{tab:cv}
  \resizebox{\textwidth}{!}{\begin{tabular}{lrrrrr}
\toprule
\textbf{Fold} & \textbf{Accuracy} & \textbf{P-macro} & \textbf{R-macro} & \textbf{F1-macro} & \textbf{F1 pond.} \\
\midrule
1 & 0.9562 & 0.9140 & 0.9144 & 0.9118 & 0.9559 \\
2 & 0.9518 & 0.9112 & 0.9080 & 0.9077 & 0.9516 \\
3 & 0.9555 & 0.9264 & 0.9160 & 0.9195 & 0.9553 \\
4 & 0.9527 & 0.9172 & 0.9123 & 0.9140 & 0.9525 \\
5 & 0.9495 & 0.9026 & 0.8923 & 0.8955 & 0.9490 \\
Media & 0.9531 & 0.9143 & 0.9086 & 0.9097 & 0.9529 \\
Desv. & 0.0027 & 0.0087 & 0.0096 & 0.0090 & 0.0028 \\
\bottomrule
\end{tabular}
}
\end{table}

\begin{figure}[H]
  \centering
  \includegraphics[width=0.82\textwidth]{cross_validation.pdf}
  \caption{Macro-F1 por fold y media de la configuración seleccionada.}
  \label{fig:cv}
\end{figure}

Esta validación no vuelve ciega la selección: Optuna y la comparación ya habían
visto una división del mismo conjunto de desarrollo. Su función es comprobar
si el resultado depende de un corte afortunado. El estimador no consultado
permanece en el holdout.

Los cinco folds produjeron 0.9097 de macro-F1 medio con desviación 0.0090. El
mejor fue el tercero (0.9195) y el menor el quinto (0.8955); la distancia entre
ambos fue 0.0240. Accuracy varió mucho menos, con media 0.9531 y desviación
0.0027. Esa diferencia entre la estabilidad de accuracy y macro-F1 vuelve a
mostrar que el movimiento ocurre en las clases pequeñas, no en el grueso del
dataset.

Potasio fue al mismo tiempo la clase con menor F1 medio (0.7039) y la más
variable (desviación 0.0496). En el quinto fold su recall cayó a 0.5524 y su F1
a 0.6304, aun cuando el accuracy global permaneció en 0.9495. Nitrógeno quedó
segundo entre los promedios bajos, con F1 0.8326, pero su desviación fue 0.0159.
Roya común y necrosis letal superaron 0.98 de F1 medio con variaciones menores
a 0.002.

La configuración es estable para las clases grandes y visualmente distintivas,
pero no para potasio. El fold bajo no invalida el ensemble; evita resumirlo como
0.91 sin decir que una partición produjo 0.63 de F1 en el problema central de
la etapa anterior. Esa variación justifica mantener métricas por clase en el
monitoreo del prototipo.
